\documentclass[10pt,twocolumn]{article}
\usepackage[scaled]{helvet}
\renewcommand\familydefault{\sfdefault}
\usepackage[T1]{fontenc}
\usepackage[margin=0.7in]{geometry}
\usepackage{titlesec}
\usepackage{enumitem}
\usepackage{parskip}
\setlength{\columnsep}{0.24in}
\setlist[itemize]{leftmargin=1.2em, topsep=0.2em, itemsep=0.18em}
\titleformat{\section}{\large\bfseries}{\thesection}{1em}{}

\begin{document}
\title{\vspace{-1.6cm}Project: Model Deployment Pipeline}
\author{AWS ML Study Notes}
\date{}
\maketitle
\small

\section*{Overview}
A model deployment pipeline is the automation path that takes an approved model artifact and promotes it safely into a serving environment. It is where ML engineering becomes operational discipline.

The goal is not only speed. The goal is repeatability, traceability, rollback capability, and separation between experimentation and production mutation.

\section*{Core Concepts}
\begin{itemize}
\item Typical building blocks include a model registry or artifact store, automated validation, security checks, deployment targets, approval gates, and post deployment monitoring hooks.
\item The pipeline should bind together model version, container image version, infrastructure configuration, and release metadata so the production state is fully reconstructable.
\item Release patterns may include staging, canary, shadow, blue green, or weighted rollout depending on the serving platform and risk tolerance.
\end{itemize}

\section*{How It Works}
\begin{itemize}
\item A training or evaluation workflow registers a candidate model and stores its metrics and artifacts in a form the deployment system can trust.
\item Deployment automation validates the artifact, updates the serving configuration, shifts traffic according to the selected release strategy, and watches health and quality signals.
\item If checks fail, the pipeline rolls back automatically or halts for investigation, preserving both uptime and release confidence.
\end{itemize}

\section*{AI and ML Relevance}
\begin{itemize}
\item This project is central to production ML because a good model in a notebook has no value unless it can be released safely and repeatedly.
\item It also enables governance by making promotion criteria explicit instead of relying on ad hoc approvals in chat or email.
\end{itemize}

\section*{Operational Notes}
\begin{itemize}
\item Do not deploy from loose files or untracked notebook outputs. Promote only immutable, versioned artifacts produced by trusted automation.
\item Separate technical health checks from model quality checks. A service can be up while prediction behavior is wrong.
\item Rollback design should be exercised before an incident. The first rollback should not happen under panic conditions.
\end{itemize}

\section*{What To Remember}
\begin{itemize}
\item The deployment pipeline is the contract between data science output and production reality.
\item Immutable artifacts and observable rollout strategy are the backbone of safe release.
\item If production cannot be traced back to a specific model version and image version, the pipeline is incomplete.
\end{itemize}

\end{document}
